% ----------------------------------------------------------------------
% Capítulo 3: Procedimentos metodológicos
% ----------------------------------------------------------------------
\chapter{Procedimentos metodológicos}
\label{cap:metodologia}

\instrucao{De quatro a oito páginas. Classifique a pesquisa (abordagem, natureza, objetivos, procedimentos), descreva o campo, os participantes, os instrumentos e as etapas de elaboração e validação do produto. Escreva no passado.}

\section{Caracterização da pesquisa}
A pesquisa é aplicada quanto à natureza, qualitativa quanto à abordagem e descritiva quanto aos objetivos. Os procedimentos combinaram análise documental, levantamento por questionário e oficinas de validação.

\section{Campo e participantes}
\instrucao{Informe quantos locais e quantas pessoas participaram, como foram escolhidos e o parecer do comitê de ética, quando houver.}

Participaram as equipes de doze bibliotecas escolares, escolhidas por conveniência entre as unidades com profissional fixo. A pesquisa seguiu a Resolução CNS nº 510/2016 \cite{brasil2016}.

\section{Etapas}
\instrucao{Descreva as etapas na ordem em que aconteceram, com o produto de cada uma.}
\begin{enumerate}[label=\arabic*)]
  \item análise de documentos das escolas (regimentos, inventários e termos de doação);
  \item aplicação do questionário do Apêndice~\ref{ap:questionario};
  \item redação da primeira versão do guia;
  \item duas oficinas de validação, com registro das sugestões;
  \item redação da versão final, avaliada com a ficha do Anexo~\ref{an:ficha}.
\end{enumerate}

\section{Análise dos dados}
As respostas abertas foram agrupadas em categorias definidas a partir dos elementos do Quadro~\ref{qua:elementos}. As respostas fechadas foram resumidas em frequências absolutas e percentuais.
